
suite les immersions \(x_i:S\to C\) (\(i=0,1\)) sont régulières, et restent des immersions régulières \(x_i:s\to C_s\) par changement de base \(s\to S\). Notons que en vertu de la suite exacte d'homotopie (IX 1.1), \(E_t\in K_\bullet(X_t\times_{k(t)}C_t)\) provient d'un élément \(E\in K_\bullet(X\times_SC)\), donc l'élément \(E_t(x_{1t})-E_t(x_{0t})\) de \(K_\bullet(X_t)\) provient de l'élément \(E(x_1)-E(x_0)\) de \(K_\bullet(X\times_SC)\), et par suite son image dans \(K_\bullet(X_s)\) par (7.3.1) n'est autre que
\[
  i^*(E(x_1)-E(x_0)).
  \tag{*}
\]
Bien entendu, pour une section \(x\) de \(C\) sur \(S\), on pose
\[
  E(x)=i_x^*(E),\qquad\text{où } i_x:X\longrightarrow X\times_SC
\]
est défini comme plus haut, et peut être considéré comme déduit de \(x\) par le changement de base plat \(X\times_SC\to C\); c'est donc une immersion régulière, à fortiori un morphisme parfait, ce qui donne bien un sens à l'expression \(E(x)\). Pour la même raison l'expression \(E_s(x_s)\) est définie, et par transitivité des images inverses (IV 2.12) dans le diagramme commutatif de morphismes parfaits suivant :
\[
\begin{tikzcd}[column sep=large,row sep=large]
X \arrow[r,"i_x"] & X\times_SC\\
X_s \arrow[u,"i"] \arrow[r,"i_{x_s}"] & X_s\times_{k(s)}C_s \arrow[u]
\end{tikzcd}
\;,
\]
on trouve que le premier membre de \((*)\) est égal à
\[
  E_s(x_{1s})-E_s(x_{0s}).
  \tag{**}
\]
Comme \(k(s)\) est algébriquement clos, les composantes irréductibles réduites de \(C_s\) sont géométriquement intègres sur \(k(s)\) et leurs points d'intersections rationnels sur \(k(s)\). On en conclut aisément que l'expression \((**)\) est algébriquement équivalente à zéro dans \(K_\bullet(X_s)\), en utilisant la connexité
